\chapter*{Résumé}
\addcontentsline{toc}{chapter}{Résumé}
\markboth{Résumé}{}

\noindent\textbf{Contexte général.} Ce rapport s'inscrit dans le cadre d'un stage de
perfectionnement de deux mois effectué au sein du département Informatique (IT) du
Mövenpick Hôtel de Tanger. Le réseau local de l'hôtel voit se connecter, au fil du
temps, un grand nombre d'équipements (postes fixes, imprimantes, terminaux mobiles,
objets connectés), sans que le département IT dispose d'un moyen simple d'en garder un
inventaire à jour.

\vspace{0.3cm}
\noindent\textbf{Objectif du projet.} L'objectif du projet est de développer un outil de
découverte et de cartographie des équipements connectés au réseau local de l'hôtel,
baptisé \textbf{NetSentry}, capable de scanner le réseau, d'identifier les équipements
actifs (adresse IP, adresse MAC, fabricant, ports ouverts), de conserver un historique
des scans, et de signaler automatiquement tout équipement inconnu apparaissant sur le
réseau.

\vspace{0.3cm}
\noindent\textbf{Méthode.} La réalisation du projet s'appuie sur une mise à niveau sur
les fondamentaux réseau (adressage IP/MAC, protocole ARP, notion de ports), puis sur la
conception et le développement d'une architecture en trois couches : un moteur de scan
reposant sur Nmap, une API backend (Node.js/Express) assurant l'orchestration, la
persistance (PostgreSQL) et la détection d'anomalies, ainsi qu'un tableau de bord web
(React) permettant à un administrateur IT de visualiser les équipements, l'historique
des scans et les alertes. La solution est conteneurisée avec Docker afin de faciliter
son déploiement.

\vspace{0.3cm}
\noindent\textbf{Résultats.} L'outil développé permet de scanner le réseau local, de
recenser les équipements connectés avec leur fabricant et leurs ports ouverts, de
conserver un historique complet des scans, et de générer une alerte dès qu'un équipement
non reconnu se connecte au réseau. Le fonctionnement de l'ensemble de la chaîne --
scan, backend, tableau de bord et déploiement conteneurisé -- a été vérifié de bout en
bout durant le stage.

\vspace{0.5cm}
\noindent \textbf{Mots-clés :} Découverte réseau, Cartographie d'équipements, Nmap,
Détection d'anomalies, Tableau de bord, Docker, NetSentry.
